% generated from papers/XXVII-nonabelian-groupoids/paper_XXVII_nonabelian_groupoids.tex -- edit the paper, then rerun assemble.py
\chapter[Non-abelian Bost--Connes groupoids]{Non-abelian Bost--Connes groupoids:\\ two no-go theorems, and a construction that realizes $\Prob(G/N_\beta)$\\ exactly when the Frobenius lifts mix}\label{ch:XXVII}
\chapterorigin{This chapter is Paper XXVII of the series (\texttt{papers/XXVII-nonabelian-groupoids/}).}
\begin{summary}
For a compact group $G$ with Frobenius conjugacy classes $\Frob_\pp$, $\pp\in S$, the Tannakian
subgroup $N_\beta=\bigcap\{\ker\rho:\Xi_\beta(\rho,S)<\infty\}$ of Chapter~\ref{ch:VI} is the canonical
non-abelian replacement of the obstruction group, and Chapter~\ref{ch:VII} showed that the obvious
groupoid model, in which the Frobenius elements act through a homomorphism of the abelian
ideal group, has $\mathrm{KMS}_\beta$ simplex $\Prob(E_\beta\backslash G)$ with $E_\beta$
abelian, which is $\Prob(G/N_\beta)$ only when $E_\beta$ is normal. We examine the two routes
proposed in Chapter~\ref{ch:VI} for closing the gap and settle both.

\emph{Noncommutative coefficient algebras} (Route B) cannot work, for a reason independent of
the representation: a $\mathrm{KMS}_\beta$ state of $\cB\rtimes J_S$ is tracial on $\cB$, so the
simplex is the simplex of scaling traces of $\cB$, inner twists are invisible, and for
$\cB=C(Y_S)\otimes M_n(\C)$ with any inner twist the equilibrium simplex is the abelian one.

\emph{Groupoids with isotropy} $N_\beta$ (the first reading of Route A) cannot work either:
by Neshveyev's theorem the isotropy contributes traces on $C^*(N_\beta)$, i.e.\ \emph{dual}
data indexed by $\Irr(N_\beta)$, and the extreme point given by the trivial representation
is fixed by every symmetry, so the symmetry group never acts transitively on the extreme
points, as it must for $\Prob(G/N_\beta)$.

\emph{Principal groupoids with a non-abelian Frobenius cocycle} (the second reading of Route
A) do work, but not canonically. Ordering the primes by norm, the ordered products
$P_n(x)=\Frob_{\pp_1}^{x_1}\cdots\Frob_{\pp_n}^{x_n}$ define a cocycle $c$ on the tail
relation with genuinely non-commuting increments, for any choice of lifts of the classes to
elements. We construct the corresponding \'etale groupoid, prove that the
$\mathrm{KMS}_\beta$ simplex of its corner algebra is $\Prob(H\backslash G)$ with $H$ the Mackey
range of $c$, and prove a \emph{non-abelian Kakutani theorem}: $\rho\circ c$ is a coboundary
if and only if $\Xi_\beta(\rho,S)=\sum_\pp\Ns\pp^{-\beta}\|I-\rho(\Frob_\pp)\|_{\HS}^2<\infty$,
the proof being a matrix martingale whose $L^2$-growth is governed by
$\det(I+t(I-U)^*(I-U)/(1-t)^2)$. Consequently $N_\beta$ is the normal closure of $H$, and
the model realizes $\Prob(G/N_\beta)$ precisely when $H$ is normal. A mixing criterion on the
products of the means $\E[\rho(\Frob_\pp)^{x_\pp}]$ forces $H=G$; for $G=S_3$ with the
transposition class at every prime, the constant lift $(12)$ gives three
$\mathrm{KMS}_\beta$ states while the alternating lift $(12),(13),(12),\dots$ gives one. The
simplex therefore depends on the lift, the lift-independent content is exactly $N_\beta$,
and there is no canonical non-abelian Bost--Connes groupoid; but whenever $N_\beta$ has no
proper subgroup with normal closure $N_\beta$ --- for instance $N_\beta=A_3$, realized by
$3$-cycle classes on a divergent set and transposition classes on a convergent one --- the
simplex is $\Prob(G/N_\beta)$ for every lift.
\end{summary}

\section{The question, and what is already known}

Throughout, $K$ is a number field, $S$ an infinite set of finite primes with
$\zeta_{K,S}(\beta)=\infty$, i.e.\ $\sum_{\pp\in S}\Ns\pp^{-\beta}=\infty$, listed by
non-decreasing norm as $\pp_1,\pp_2,\dots$, with $t_i=\Ns\pp_i^{-\beta}\in(0,t_{\max}]$,
$t_{\max}=2^{-\beta}<1$. We write
$V=\N^S$ for the space of valuation vectors with the product measure
$\pi_\beta=\bigotimes_i\mu_{t_i}$, $\mu_t(k)=(1-t)t^k$, forced on every $\mathrm{KMS}_\beta$
state by [Ch.~\ref{ch:Main}, Lem.~2.10]; $\cR_S$ is the tail equivalence relation on $V$ (finitely
many coordinates changed), ergodic by the zero--one law [Ch.~\ref{ch:Main}, Lem.~3.1];
$\cF_{\le n}=\sigma(x_1,\dots,x_n)$ and $\cF_{>n}=\sigma(x_{n+1},x_{n+2},\dots)$.

$G$ is a compact second countable group; for each $i$ a conjugacy class $C_i\subset G$ is
given (the Frobenius class at $\pp_i$, or a Satake class), and a \emph{lift} is a choice of
elements $u_i\in C_i$. For a unitary representation $\rho$ of $G$ put
\begin{equation}\label{XXVII:eq:Xi}
\Xi_\beta(\rho,S)=\sum_i t_i\,\bigl\|I-\rho(u_i)\bigr\|_{\HS}^2
=2\sum_it_i\bigl(\dim\rho-\Re\Tr\rho(u_i)\bigr),
\end{equation}
which depends only on the classes, and let $\cT_\beta=\{\rho:\Xi_\beta(\rho,S)<\infty\}$ and
$N_\beta=\bigcap_{\rho\in\cT_\beta}\ker\rho$; $\cT_\beta=\Rep(G/N_\beta)$ is the Tannakian
theorem of Chapter~\ref{ch:VI}, and $N_\beta$ depends only on the classes.

The abelian theory has three ingredients: a \emph{space} on which the ideal semigroup acts,
whose $\mathrm{KMS}_\beta$ measures are the equilibrium states; a \emph{symmetry group}
permuting them; and a \emph{cocycle} whose coboundary question is the Kakutani dichotomy.
The simplex is $\Prob(G_S/\Xi_\beta^\perp)$ with $G_S$ acting transitively on the extreme
points [Ch.~\ref{ch:Main}, Thm.~3.6]. The target of a non-abelian theory is a system whose
$\mathrm{KMS}_\beta$ simplex is $\Prob(G/N_\beta)$ with $G$ acting transitively on the
extreme points by translation, stabilizer $N_\beta$. The model of Chapter~\ref{ch:VII} --- $I_S$
acting on $V\times G$ through the homomorphism $a\mapsto\prod u_i^{a_i}$, which requires
commuting lifts --- gives $\Prob(E_\beta\backslash G)$ with $E_\beta$ contained in the closed
abelian group generated by the $u_i$, and $N_\beta$ is the normal closure of $E_\beta$.
Two ways around this were proposed in Chapter~\ref{ch:VI} (its Remark 6.1) and left open in Chapter~\ref{ch:VII}:
a groupoid built directly, whose unit space carries the tail relation and a $G$-action and
whose isotropy realizes $N_\beta$; and a noncommutative coefficient algebra carrying
$\rho$, with the commutative Ore semigroup $J_S$ kept. We take the second first, because
its failure is the simplest.

\section{Route B: noncommutative coefficient algebras}\label{XXVII:sec:B}

Let $\cB$ be a unital $C^*$-algebra with an action $\alpha$ of $J_S=\N^{(S)}$ by injective
endomorphisms with hereditary ranges, $\alpha_\aaa(1)$ a projection, so that Laca's dilation
applies \cite{Laca}: $\cB\rtimes_\alpha J_S$ is the full corner $\mathbf 1_\cB(\widetilde\cB\rtimes I_S)\mathbf 1_\cB$
of the crossed product of the direct limit $\widetilde\cB=\varinjlim(\cB,\alpha)$ by the
automorphic extension of $\alpha$ to the group $I_S$; it is generated by $\cB$ and isometries
$\mu_\aaa$ with $\mu_\aaa b\mu_\aaa^*=\alpha_\aaa(b)$, and the dynamics is
$\sigma_t(\mu_\aaa)=\Ns\aaa^{it}\mu_\aaa$, $\sigma_t|_\cB=\id$. The gauge action of the
compact group $\widehat{I_S}$ gives the faithful conditional expectation
$E:\cB\rtimes_\alpha J_S\to\cB$, the average of the gauge automorphisms; so ``a tractable
faithful expectation onto the coefficient algebra'' exists for every $\cB$. It does not
help. Write $T(\cB)$ for the tracial state space of $\cB$, a Choquet simplex, and
\[
T_\beta(\cB,\alpha)=\bigl\{\tau\in T(\cB):\ \tau\circ\alpha_\aaa=\Ns\aaa^{-\beta}\tau\ \ (\aaa\in J_S)\bigr\},
\]
the \emph{scaling traces}, a closed face of $T(\cB)$.

\begin{theorem}[Route B is abelian]\label{XXVII:thm:B}
\begin{enumerate}
\item Every $\mathrm{KMS}_\beta$ state $\varphi$ of $\cB\rtimes_\alpha J_S$ restricts to a
scaling trace $\tau=\varphi|_\cB\in T_\beta(\cB,\alpha)$.
\item For every $\tau\in T_\beta(\cB,\alpha)$, $\tau\circ E$ is a $\mathrm{KMS}_\beta$ state
with restriction $\tau$; so $\tau\mapsto\tau\circ E$ is an affine injection
$T_\beta(\cB,\alpha)\to\mathrm{KMS}_\beta(\cB\rtimes_\alpha J_S)$.
\item If the numbers $\log\Ns\pp$, $\pp\in S$, are linearly independent over $\Q$ --- for
instance $K=\Q$ --- then every $\mathrm{KMS}_\beta$ state is of the form $\tau\circ E$, and
$\mathrm{KMS}_\beta(\cB\rtimes_\alpha J_S)\cong T_\beta(\cB,\alpha)$ affinely.
\end{enumerate}
Consequently: (a) if $\alpha'_\aaa=\mathrm{Ad}(w_\aaa)\circ\alpha_\aaa$ with unitaries
$w_\aaa\in\cB$ is another action, then $T_\beta(\cB,\alpha')=T_\beta(\cB,\alpha)$; inner twists
are invisible. (b) For $\cB=C(Y_S)\otimes M_n(\C)$ and any action of the form
$\alpha_\aaa\otimes\mathrm{Ad}(v_\aaa)$, $v_\aaa\in U(n)$, the scaling traces are
$\nu\otimes\frac1n\Tr$ with $\nu$ a $\mathrm{KMS}_\beta$ measure of the classical system, so
$T_\beta=\Prob(G_S/\Xi_\beta^\perp)$ [Ch.~\ref{ch:Main}, Thm.~3.6]: the abelian answer. (c) In every
case the equilibrium theory of $\cB\rtimes_\alpha J_S$ is that of the affine action of
$J_S$ on the Choquet simplex $T(\cB)$; the noncommutativity of $\cB$ is quotiented out by
traciality before any symmetry can act on it.
\end{theorem}

\begin{proof}
(1) $\cB$ lies in the fixed-point algebra of $\sigma$, and for $x,y$ there the KMS condition
gives $\varphi(xy)=\varphi(y\sigma_{i\beta}(x))=\varphi(yx)$; so $\tau$ is tracial. For
$b\in\cB$, the KMS condition with the entire element $\mu_\aaa$,
$\sigma_{i\beta}(\mu_\aaa)=\Ns\aaa^{-\beta}\mu_\aaa$, gives
$\tau(\alpha_\aaa(b))=\varphi(\mu_\aaa b\mu_\aaa^*)=\Ns\aaa^{-\beta}\varphi(b\mu_\aaa^*\mu_\aaa)=\Ns\aaa^{-\beta}\tau(b)$.

(2) Extend $\tau$ to the semifinite trace $\tilde\tau$ on $\widetilde\cB$ determined by
$\tilde\tau(\alpha_\aaa^{-1}(b))=\Ns\aaa^{\beta}\tau(b)$, consistent by the scaling relation,
and let $\gamma$ be the automorphic action of $I_S$ on $\widetilde\cB$, so that
$\tilde\tau\circ\gamma_a=\Ns\aaa^{-\beta}\tilde\tau$ for $a\in I_S$. On the dense span of
elements $x=a_1u_a$, $y=a_2u_b$ of $\widetilde\cB\rtimes I_S$ ($a_i\in\widetilde\cB$ in the
domain of $\tilde\tau$, $u_a$ the implementing unitaries),
\[
(\tilde\tau\circ E)(xy)=\delta_{a+b,0}\,\tilde\tau\bigl(a_1\gamma_a(a_2)\bigr),\qquad
(\tilde\tau\circ E)\bigl(y\,\sigma_{i\beta}(x)\bigr)=\delta_{a+b,0}\,\Ns\aaa^{-\beta}\tilde\tau\bigl(a_2\gamma_{-a}(a_1)\bigr),
\]
and $\Ns\aaa^{-\beta}\tilde\tau(a_2\gamma_{-a}(a_1))=\Ns\aaa^{-\beta}\Ns\aaa^{\beta}\tilde\tau(\gamma_a(a_2)a_1)=\tilde\tau(a_1\gamma_a(a_2))$,
using $\tilde\tau\circ\gamma_{a}=\Ns\aaa^{-\beta}\tilde\tau$ read as $\tilde\tau=\Ns\aaa^{\beta}\tilde\tau\circ\gamma_a$
and then traciality. So $\tilde\tau\circ E$ satisfies the KMS condition on the corner,
where it is the state $\tau\circ E$. Injectivity of $\tau\mapsto\tau\circ E$ is clear.

(3) A $\mathrm{KMS}_\beta$ state is $\sigma$-invariant. Under the hypothesis, the one-parameter
group $t\mapsto(\Ns\pp^{it})_\pp$ is dense in $\widehat{I_S}=\prod_\pp\T$ by Kronecker's
theorem applied to finite subsets of $S$, so a $\sigma$-invariant state is invariant under
the whole gauge action and hence factors through $E$: $\varphi=\varphi\circ E=\tau\circ E$
with $\tau=\varphi|_\cB$, which is a scaling trace by (1).

(a) A trace is invariant under inner automorphisms, so $\tau\circ\alpha'_\aaa=\tau\circ\alpha_\aaa$.
(b) A tracial state of $C(Y_S)\otimes M_n$ is $\nu\otimes\frac1n\Tr$, and
$\Tr(v_\aaa m v_\aaa^*)=\Tr m$ reduces the scaling condition to $\nu\circ\alpha_\aaa=\Ns\aaa^{-\beta}\nu$.
(c) is (1)--(3).
\end{proof}

\begin{remark}
Without the independence hypothesis in (3), $\sigma$-invariance only forces $\varphi$ to
factor through the expectation onto the fixed-point algebra of $\sigma$, which also
contains the elements $\mu_\aaa b\mu_\bb^*$ with $\Ns\aaa=\Ns\bb$, $\aaa\ne\bb$; for the
classical system these carry nothing because the groupoid is principal on a conull set
[Ch.~\ref{ch:Main}, Prop.~2.8], and for general $\cB$ there may be additional $\mathrm{KMS}_\beta$
states, all of them tracial on the fixed-point algebra. None of this reintroduces the
representation: the matrix Kakutani dichotomy of Chapter~\ref{ch:VI}, valid for commuting lifts,
would enter Route B through the twist $\mathrm{Ad}\,\rho(u_\pp)$, and (a) says the twist is
never seen; for non-commuting lifts $\alpha_\pp\otimes\mathrm{Ad}\,\rho(u_\pp)^{-1}$ do not
even commute, so there is no action to twist. The modular group is not ``broken'' by a
noncommutative $\cB$: it is the same gauge dynamics, and it forces the equilibrium state
to forget everything about $\cB$ except its traces. Route B is closed.
\end{remark}

\section{Route A, first reading: isotropy cannot carry $N_\beta$}\label{XXVII:sec:iso}

Let $\cG$ be a second countable locally compact Hausdorff \'etale groupoid with unit space
$\Omega$, and $c:\cG\to\R$ a continuous cocycle defining $\sigma_t(f)(\gamma)=e^{itc(\gamma)}f(\gamma)$
on $C^*(\cG)$ (the Haar system is counting measure on the fibres, as for every \'etale
groupoid). Neshveyev's theorem \cite[Thm.~1.3]{Nesh13} describes the $\mathrm{KMS}_\beta$
states: they are the functionals
\begin{equation}\label{XXVII:eq:nesh}
\varphi(f)=\int_\Omega\sum_{\gamma\in\cG_x^x}f(\gamma)\,\phi_x(u_\gamma)\,d\mu(x),
\end{equation}
where $\mu$ is a quasi-invariant probability measure with Radon--Nikodym cocycle $e^{-\beta c}$
and $x\mapsto\phi_x$ is a $\mu$-measurable field of tracial states on the group
$C^*$-algebras $C^*(\cG_x^x)$ of the isotropy groups, vanishing on $u_\gamma$ unless
$c(\gamma)=0$, and conjugation-equivariant: $\phi_{r(\gamma)}(u_{\gamma h\gamma^{-1}})=\phi_{s(\gamma)}(u_h)$
for $\mu$-a.e.\ $\gamma$. The extreme points are the pairs with $\mu$ ergodic and $\phi$
extremal among equivariant fields.

\begin{theorem}[No-go for isotropy]\label{XXVII:thm:isoNoGo}
Let $(\cG,c)$ be as above, and let a group $G$ act on $\cG$ by automorphisms preserving $c$,
hence on $\mathrm{KMS}_\beta(C^*(\cG))$ by affine bijections. For an extreme
$\mathrm{KMS}_\beta$ state with quasi-invariant measure $\mu$ let $\mathcal E_\mu$ be the set
of extreme $\mathrm{KMS}_\beta$ states with measure $\mu$.
\begin{enumerate}
\item The state $\varphi_\mu(f)=\int_\Omega f|_\Omega\,d\mu$ --- the field $\phi_x=$ trivial
character --- belongs to $\mathcal E_\mu$, and it is fixed by the stabilizer $G_\mu$ of
$\mu$ in $G$.
\item If the isotropy groups $\cG_x^x$ are $\mu$-a.e.\ isomorphic to a fixed compact group
$N\subseteq\ker c$, with the isotropy bundle constant and the conjugation action of $\cG$
on it trivial, then $\mathcal E_\mu$ is in bijection with $\Irr(N)$, via
$\rho\mapsto(\mu,\phi_x\equiv\chi_\rho/\dim\rho)$.
\item Consequently, if $G$ acts transitively on the extreme $\mathrm{KMS}_\beta$ states, then
$\mathcal E_\mu=\{\varphi_\mu\}$ for every $\mu$, and in the situation of (2) $N=\{1\}$: the
isotropy is $\mu$-a.e.\ trivial.
\end{enumerate}
In particular no groupoid whose $\mathrm{KMS}_\beta$ multiplicity is produced by isotropy
has a simplex of the form $\Prob(G/N)$ with $G$ acting transitively on the extreme points
and $N\ne1$.
\end{theorem}

\begin{proof}
(1) The trivial character is a tracial state of every $C^*(\cG_x^x)$, vanishes on $u_\gamma$
for $c(\gamma)\ne0$ as required, and the constant field is conjugation-equivariant; so
$\varphi_\mu$ is a $\mathrm{KMS}_\beta$ state by \eqref{XXVII:eq:nesh}, and it is extreme because
$\mu$ is ergodic and the trivial character is an extreme tracial state. An automorphism of
$(\cG,c)$ transforms a pair $(\mu,\phi)$ into $(\mu\circ g^{-1},\phi\circ g^{-1})$; if it
fixes $\mu$ it fixes the pair with the trivial field.

(2) Tracial states of $C^*(N)$ are the convex combinations of normalized characters
$\chi_\rho/\dim\rho$, $\rho\in\Irr(N)$, the extreme ones being the normalized characters.
With constant bundle and trivial conjugation action, equivariance says the field is
$\cG$-invariant, hence a.e.\ constant by ergodicity of $\mu$; so the equivariant fields
are the constant ones, the extreme ones the constant normalized characters.

(3) The map $\varphi\mapsto\mu$ from extreme $\mathrm{KMS}_\beta$ states to ergodic
quasi-invariant measures is $G$-equivariant, so if $G$ is transitive on the extreme states
then $G_\mu$ is transitive on the fibre $\mathcal E_\mu$; by (1) it fixes $\varphi_\mu\in\mathcal E_\mu$,
so $\mathcal E_\mu=\{\varphi_\mu\}$, and by (2) $|\Irr(N)|=1$, i.e.\ $N=1$. For the last
sentence, the translation action of $G$ on $\Prob(G/N)$ is transitive on the extreme
points.
\end{proof}

\begin{remark}
The theorem says that isotropy is the wrong place to put $N_\beta$: it produces the
\emph{dual} object $\Irr(N_\beta)$, with the trivial representation as a distinguished,
symmetry-fixed equilibrium state, whereas symmetry breaking needs a homogeneous space
with no distinguished point. ``Isotropy $=N_\beta$'' would give a simplex with a canonical
symmetric state and $|\Irr(N_\beta)|$ extreme points over each ergodic measure, which is
$\Prob(G/N_\beta)$ neither as a $G$-space nor, in general, as a set. The classical system
itself illustrates the point: its groupoid is principal on a conull set
[Ch.~\ref{ch:Main}, Prop.~2.8], and all of its multiplicity comes from the ergodic decomposition
of the unit space. That is where a non-abelian construction must put it.
\end{remark}

\section{Route A, second reading: the ordered Frobenius cocycle}\label{XXVII:sec:ordered}

\subsection{The cocycle}

Fix lifts $u_i\in C_i$. For $x\in V$ and $n\ge0$ put
\begin{equation}\label{XXVII:eq:Pn}
P_n(x)=u_1^{x_1}u_2^{x_2}\cdots u_n^{x_n}\in G,
\end{equation}
an ordered product, and for $(y,x)\in\cR_S$ let $n$ be any index beyond the last coordinate
in which $x$ and $y$ differ and
\begin{equation}\label{XXVII:eq:c}
c(y,x)=P_n(y)\,P_n(x)^{-1}.
\end{equation}
This is independent of $n$: $P_{n'}(x)=P_n(x)u_{n+1}^{x_{n+1}}\cdots u_{n'}^{x_{n'}}$ and the
extra factors agree for $x$ and $y$ and cancel on the right. It is a cocycle,
$c(z,y)c(y,x)=c(z,x)$, being of the form $P(y)P(x)^{-1}$ on each $\cR_{F_n}$. Its
increments are
\begin{equation}\label{XXVII:eq:incr}
c(x+e_j,x)=P_{j-1}(x)\,u_j\,P_{j-1}(x)^{-1},
\end{equation}
the $j$-th lift conjugated by the ordered product of the earlier ones; for commuting lifts
this is $u_j$ and $c$ is the constant-increment cocycle of Chapter~\ref{ch:VII}, but in general the
increments at different coordinates do not commute, and \eqref{XXVII:eq:c} is the way to have a
cocycle nonetheless: the obstruction of Chapter~\ref{ch:VI}, that a cocycle with \emph{constant}
non-commuting increments cannot exist on a relation generated by commuting shifts, is
circumvented by letting the increment at $\pp_j$ depend on the coordinates before it.

\subsection{The groupoid and its algebra}

To have a topological model we compactify each coordinate together with its lift. Let
\[
X_i=\overline{\{(k,u_i^k):k\in\N\}}\subseteq\overline\N\times G,\qquad X=\prod_{i\ge1}X_i,
\]
compact, second countable, with $\N\subset X_i$ open and dense (as $\{k\}\times\{u_i^k\}$) and
boundary $\partial X_i=\{\infty\}\times\overline{\{u_i^k\}}$. Write $g_i:X_i\to G$ for the
second projection, continuous, with $g_i(k)=u_i^k$. The shift $k\mapsto k+1$ extends to a
homeomorphism $s_i$ of $X_i$ onto $X_i\setminus\{0\}$, $(\infty,g)\mapsto(\infty,u_ig)$, and
$J_S$ acts on $X$ by $x\mapsto x+a$ coordinatewise, injectively, with clopen images. Put
\[
P_n(x)=g_1(x_1)\cdots g_n(x_n),\qquad
c(x+e_j,x)=P_{j-1}(x)u_jP_{j-1}(x)^{-1},
\]
continuous on $X$ and extending \eqref{XXVII:eq:Pn}--\eqref{XXVII:eq:incr}. Let $\widetilde X=\varinjlim(X,+a)$
be the dilation, a locally compact space on which $I_S=\Z^{(S)}$ acts by homeomorphisms, and
extend $c$ to $\widetilde X$ by the same formulas. Define the action of $I_S$ on
$\Omega=\widetilde X\times G$ by
\begin{equation}\label{XXVII:eq:action}
a\cdot(x,g)=\bigl(x+a,\ c(x+a,x)\,g\bigr),
\end{equation}
which is an action because $c$ is a cocycle, and by homeomorphisms because $c$ is
continuous. Let
\[
\cG_c=I_S\ltimes\Omega,\qquad
\cA_c=\mathbf 1_{X\times G}\,C^*(\cG_c)\,\mathbf 1_{X\times G}=C(X\times G)\rtimes J_S,
\]
the transformation groupoid with counting measures as Haar system \cite{Renault}, and
its unital full corner; the dynamics is $\sigma_t(\mu_\aaa)=\Ns\aaa^{it}\mu_\aaa$, given by the continuous
cocycle $(a,\omega)\mapsto\log\Ns\aaa$ on $\cG_c$. The group $G$ acts on $\Omega$ by right
translation $R_h(x,g)=(x,gh)$, commuting with \eqref{XXVII:eq:action}, hence on $\cA_c$ by
$\sigma$-commuting automorphisms.

\begin{lemma}\label{XXVII:lem:basic}
$\cG_c$ is second countable, \'etale, Hausdorff and amenable, principal on the conull set
$V\times G$ (points with all coordinates finite), and $\cA_c$ is nuclear. The
$\mathrm{KMS}_\beta$ states of $\cA_c$ are in affine bijection with the probability
measures $\nu$ on $X\times G$ with $\nu(a\cdot E)=\Ns\aaa^{-\beta}\nu(E)$ for $a\in J_S$; every
such $\nu$ is concentrated on $V\times G$, has $V$-marginal $\pi_\beta$, and disintegrates as
$\nu=\pi_\beta\otimes\rho_x$ with a measurable family $\rho_x\in\Prob(G)$ satisfying
\begin{equation}\label{XXVII:eq:equiv}
\rho_{x+e_j}=\bigl(c(x+e_j,x)\bigr)_*\rho_x\qquad\text{for }\pi_\beta\text{-a.e.\ }x,\ \text{every }j .
\end{equation}
\end{lemma}

\begin{proof}
$I_S$ is abelian, so $\cG_c$ is amenable and $C^*=C^*_r$; \'etaleness and Hausdorffness are
those of a transformation groupoid by homeomorphisms. If $a\cdot(x,g)=(x,g)$ then $x+a=x$,
which forces $a=0$ when all $x_i$ are finite; so the isotropy is trivial on $V\times G$. For
the KMS states, the argument of [Ch.~\ref{ch:Main}, Prop.~2.8] applies verbatim (or one may invoke
\eqref{XXVII:eq:nesh}): a $\mathrm{KMS}_\beta$ state is $\nu\circ E$ for a scaling measure $\nu$
on the unit space of the corner. If $E\subseteq\{x_i=\infty\}$ then $e_i\cdot E=E$ and the
scaling gives $\nu(E)=t_i\nu(E)$, so $\nu(E)=0$: $\nu$ is concentrated on $V\times G$. The
$V$-marginal satisfies the scaling relations of the valuation coordinates alone, so it is
$\pi_\beta$ by [Ch.~\ref{ch:Main}, Lem.~2.10]. Disintegrating and comparing the two sides of the
scaling identity on sets $B\times C$, using $\pi_\beta(B+e_j)=t_j\pi_\beta(B)$, gives
\eqref{XXVII:eq:equiv}.
\end{proof}

\subsection{Classification by the Mackey range}

Recall Zimmer's theory of cocycles into compact groups \cite[Ch.~4 and 9]{Zimmer}, in the
setting of countable equivalence relations \cite{FM}: for an ergodic countable equivalence
relation $(V,\pi,\cR)$ and a measurable cocycle
$c:\cR\to G$ into a compact group, there is a closed subgroup $H\le G$, unique up to
conjugacy, and a measurable $f:V\to G$ such that $c'(y,x)=f(y)^{-1}c(y,x)f(x)$ takes values in
$H$ and the skew product relation $\cR\ltimes_{c'}H$ on $V\times H$ is ergodic; $H$ is the
\emph{Mackey range} (algebraic hull) of $c$, and $c$ is cohomologous to a cocycle into a
closed subgroup $H'$ iff $H'$ contains a conjugate of $H$.

\begin{theorem}[The simplex of the ordered model]\label{XXVII:thm:simplex}
Let $H$ be the Mackey range of $c$ on $(V,\pi_\beta,\cR_S)$. Then
\[
\mathrm{KMS}_\beta(\cA_c)\ \cong_{\mathrm{affine}}\ \Prob\bigl(H\backslash G\bigr),
\]
the extreme points being the families $\rho_x=$ Haar measure on $f(x)Hk$, $k\in G$, and the
right translation action of $G$ is transitive on the extreme points with stabilizer of the
point $Hk$ equal to $k^{-1}Hk$.
\end{theorem}

\begin{proof}
By Lemma~\ref{XXVII:lem:basic} we classify the equivariant families \eqref{XXVII:eq:equiv}, i.e.\ the
probability measures $\nu=\pi_\beta\otimes\rho_x$ on $V\times G$ invariant in measure class
and with the prescribed Radon--Nikodym derivative under the skew product relation
$\widetilde\cR=\cR_S\ltimes_cG$, $(x,g)\sim(y,c(y,x)g)$; the extreme ones are the ergodic
ones. The map $\Psi(x,g)=(x,f(x)^{-1}g)$ carries $\widetilde\cR$ to $\cR_S\ltimes_{c'}G$ and
preserves $\pi_\beta\otimes m_G$. The latter relation preserves each set $V\times Hk$, on which
it is isomorphic to $\cR_S\ltimes_{c'}H$ via $h\mapsto hk$, ergodic by the definition of the
Mackey range. Hence the ergodic components of $\widetilde\cR$ are the sets
$\Psi^{-1}(V\times Hk)=\{(x,g):f(x)^{-1}g\in Hk\}$, $Hk\in H\backslash G$. Within a
component, Haar on the fibres is the \emph{only} equivariant family: if
$\rho'_x\in\Prob(H)$ satisfies $\rho'_y=c'(y,x)_*\rho'_x$, then for a nontrivial
$\sigma\in\Irr(H)$ the Fourier coefficient $B(x)=\bigl(\int_H\sigma(h)^*d\rho'_x(h)\bigr)^*$
satisfies $B(y)=\sigma(c'(y,x))B(x)$, so $\Phi(x,h)=\Tr\bigl(\sigma(h)^*B(x)\bigr)$ is
$\cR_S\ltimes_{c'}H$-invariant --- $\Phi(y,c'h)=\Tr(\sigma(h)^*\sigma(c')^*\sigma(c')B(x))=\Phi(x,h)$
--- hence a.e.\ constant by ergodicity; since $\int_H\sigma(h)\,dh=0$ the constant is $0$,
and Fourier inversion in $h$ gives $B=0$. All nontrivial coefficients of $\rho'_x$ vanish,
so $\rho'_x=m_H$ a.e. Thus the ergodic
equivariant families are the Haar measures on the fibres $f(x)Hk$, and every equivariant
family is a probability mixture of these, uniquely, since distinct ergodic components carry
mutually singular measures. Right translation by $h$ maps the component of $Hk$ to that of
$Hkh$.
\end{proof}

For commuting lifts $f$ can be taken trivial and $H=E_\beta$, recovering Chapter~\ref{ch:VII}. The
content of the model is therefore the conjugacy class of $H$. Two questions remain: how
$H$ is related to $N_\beta$, and how $H$ depends on the lifts.

\subsection{The non-abelian Kakutani theorem}

For a unitary representation $\rho:G\to U(d)$ write $U_i=\rho(u_i)$ and
\begin{equation}\label{XXVII:eq:mean}
m_i=\E_{\mu_{t_i}}\bigl[U_i^{x_i}\bigr]=\sum_{k\ge0}(1-t_i)t_i^kU_i^k=(1-t_i)(I-t_iU_i)^{-1},
\end{equation}
invertible since $\|t_iU_i\|<1$, with
\begin{equation}\label{XXVII:eq:minv}
m_i^{-*}m_i^{-1}=\frac{(I-t_iU_i^*)(I-t_iU_i)}{(1-t_i)^2}=I+\frac{t_i}{(1-t_i)^2}(I-U_i)^*(I-U_i)\ \ge\ I .
\end{equation}
If $e^{i\phi_{i,1}},\dots,e^{i\phi_{i,d}}$ are the eigenvalues of $U_i$, then
\begin{equation}\label{XXVII:eq:det}
\det\bigl(m_i^{-*}m_i^{-1}\bigr)=\prod_{r=1}^d\Bigl(1+\frac{t_i|1-e^{i\phi_{i,r}}|^2}{(1-t_i)^2}\Bigr),
\qquad
\sum_r|1-e^{i\phi_{i,r}}|^2=\|I-U_i\|_{\HS}^2 .
\end{equation}
Put $M_{(n,m]}=m_{n+1}m_{n+2}\cdots m_m$ and $M_n=M_{(0,n]}$. A \emph{coboundary} means a
measurable $F:V\to U(d)$ with $\rho(c(y,x))=F(y)F(x)^{-1}$ for $\pi_\beta$-a.e.\ $(y,x)\in\cR_S$.

\begin{theorem}[Non-abelian Kakutani]\label{XXVII:thm:kakutani}
For every lift and every unitary representation $\rho$,
\[
\rho\circ c\ \text{is a coboundary}\iff\Xi_\beta(\rho,S)<\infty .
\]
\end{theorem}

\begin{proof}
Two structural facts are used throughout. First, for $F:V\to M_d(\C)$ measurable with
$F(y)=\rho(c(y,x))F(x)$ on $\cR_S$, the function $K_n=P_n^{-1}F$ (with $P_n=\rho(P_n(x))$, we
suppress $\rho$) is invariant under changes of the first $n$ coordinates: for $y\sim x$
differing only there, $K_n(y)=P_n(y)^{-1}\rho(c(y,x))F(x)=P_n(y)^{-1}P_n(y)P_n(x)^{-1}F(x)=K_n(x)$;
so $K_n$ is $\cF_{>n}$-measurable [Ch.~\ref{ch:Main}, Lem.~3.1] and independent of $P_n$, which is
$\cF_{\le n}$-measurable. Second, $\E[P_n]=M_n$ and $\E[P_m|\cF_{\le n}]=P_nM_{(n,m]}$, by
independence of the coordinates and the ordering of the product.

$(\Leftarrow)$ Let $\Xi_\beta(\rho,S)<\infty$ and $Z_n=P_nM_n^{-1}$. Then
$\E[Z_{n+1}|\cF_{\le n}]=P_nm_{n+1}M_{n+1}^{-1}=Z_n$: a matrix-valued martingale. Its
$L^2$-norm is deterministic, $\E\|Z_n\|_{\HS}^2=\Tr(M_n^{-1}M_n^{-*})=\|M_n^{-1}\|_{\HS}^2$,
and by \eqref{XXVII:eq:minv}, iterating $m_{i}^{-*}(\,\cdot\,)m_i^{-1}\le\lambda_i(\,\cdot\,)$ with
$\lambda_i=1+t_i\|I-U_i\|_{\op}^2/(1-t_i)^2$ from the inside out,
\[
\|M_n^{-1}\|_{\HS}^2\le d\prod_{i\le n}\lambda_i\le d\,\exp\Bigl(\frac{1}{(1-t_{\max})^2}\sum_it_i\|I-U_i\|_{\HS}^2\Bigr)<\infty ,
\]
using $t_i\le t_{\max}<1$. So $Z_n\to Z_\infty=:F$ a.s.\ and in $L^2$. Moreover
$|\det Z_n|=|\det M_n^{-1}|=\prod_{i\le n}\det(m_i^{-*}m_i^{-1})^{1/2}$ increases to a finite
limit $\ge1$ by \eqref{XXVII:eq:det} and $\Xi_\beta<\infty$, so $\det F\ne0$ a.s.\ and $F$ is
invertible. For $y\sim x$ and $n$ large, $Z_n(y)Z_n(x)^{-1}=P_n(y)P_n(x)^{-1}=\rho(c(y,x))$, and
letting $n\to\infty$ gives $F(y)=\rho(c(y,x))F(x)$. Then $F^*F$ is $\cR_S$-invariant,
$F(y)^*F(y)=F(x)^*\rho(c)^*\rho(c)F(x)=F(x)^*F(x)$, hence a.e.\ equal to a constant positive
matrix $Q^2$, and $W=FQ^{-1}$ is unitary with $\rho(c(y,x))=W(y)W(x)^{-1}$.

$(\Rightarrow)$ Let $F$ be a unitary coboundary and $C_n=\E[K_n]$. Then
$\E[F|\cF_{\le n}]=P_nC_n$, and comparing $\E[F|\cF_{\le n}]=\E[\E[F|\cF_{\le m}]\,|\,\cF_{\le n}]=P_nM_{(n,m]}C_m$
gives
\begin{equation}\label{XXVII:eq:consist}
C_n=M_{(n,m]}\,C_m\qquad(m>n).
\end{equation}
Suppose $\Xi_\beta(\rho,S)=\infty$. Then for each $n$, $|\det M_{(n,m]}^{-1}|\to\infty$ as
$m\to\infty$ by \eqref{XXVII:eq:det}, while $\|C_m\|_{\HS}\le\E\|K_m\|_{\HS}=\sqrt d$; so
$|\det C_m|=|\det M_{(n,m]}^{-1}|\,|\det C_n|$ bounded forces $\det C_n=0$ for every $n$. But
$P_nC_n=\E[F|\cF_{\le n}]\to F$ in $L^2$ by martingale convergence, $F$ being measurable with
respect to $\bigvee_n\cF_{\le n}$, so along a subsequence $\det(P_{n_k}C_{n_k})\to\det F$
a.e., and $|\det F|=1$; this contradicts $\det C_{n_k}=0$.
\end{proof}

\begin{corollary}[$N_\beta$ is the normal closure of the Mackey range]\label{XXVII:cor:normalclosure}
For every lift, with $H$ the Mackey range of $c$,
\[
\cT_\beta=\{\rho:\ \rho|_H=1\},\qquad N_\beta=\langle\langle H\rangle\rangle ,
\]
the smallest closed normal subgroup containing $H$. Consequently
$\mathrm{KMS}_\beta(\cA_c)=\Prob(G/N_\beta)$, with $G$ acting on the extreme points as on
$G/N_\beta$, if and only if $H$ is normal in $G$; in general there is a $G$-equivariant
surjection $H\backslash G\to G/N_\beta$ of extreme boundaries, a bijection iff $H\trianglelefteq G$.
\end{corollary}

\begin{proof}
$\rho\circ c$ is a coboundary iff the Mackey range of $\rho\circ c$ is trivial; and the Mackey
range of $\rho\circ c$ is $\rho(H)$, since $\rho\circ c'$ takes values in $\rho(H)$ and the
skew product by $\rho\circ c'$ on $V\times\rho(H)$ is a factor of the ergodic
$\cR_S\ltimes_{c'}H$. So by Theorem~\ref{XXVII:thm:kakutani}, $\Xi_\beta(\rho,S)<\infty$ iff
$\rho|_H=1$. The intersection of the kernels of the representations trivial on $H$ is the
closed normal subgroup generated by $H$, by Peter--Weyl applied to $G/\langle\langle H\rangle\rangle$.
The rest is Theorem~\ref{XXVII:thm:simplex}.
\end{proof}

\subsection{When is the Mackey range normal? A mixing criterion}

\begin{theorem}[Mixing forces $H=G$]\label{XXVII:thm:mixing}
Suppose that for every nontrivial irreducible $\rho$ and every $n$,
$\|M_{(n,m]}\|_{\op}\to0$ as $m\to\infty$. Then the skew product $\cR_S\ltimes_cG$ is ergodic,
$H=G$, and $\cA_c$ has a unique $\mathrm{KMS}_\beta$ state.
\end{theorem}

\begin{proof}
Let $\Phi\in L^\infty(V\times G)$ be $\widetilde\cR$-invariant and $A_\rho(x)=\int_G\Phi(x,h)\rho(h)^*dh$
its Fourier coefficients. Invariance gives $A_\rho(y)=A_\rho(x)\rho(c(y,x))^*$, so $B=A_\rho^*$
satisfies $B(y)=\rho(c(y,x))B(x)$ and is bounded. As in the proof of
Theorem~\ref{XXVII:thm:kakutani}, $\E[B|\cF_{\le n}]=P_nC_n$ with $C_n=\E[P_n^{-1}B]$ satisfying
\eqref{XXVII:eq:consist} and $\|C_m\|_{\HS}\le\sqrt d\,\|B\|_\infty$. By hypothesis
$\|C_n\|\le\|M_{(n,m]}\|_{\op}\sup_m\|C_m\|\to0$, so $C_n=0$ for all $n$, hence $B=0$ by
martingale convergence: $A_\rho=0$ for every nontrivial irreducible $\rho$. Thus $\Phi(x,g)$
does not depend on $g$, and then it is $\cR_S$-invariant, hence constant. Ergodicity of the
skew product means $H=G$, and Theorem~\ref{XXVII:thm:simplex} gives the rest.
\end{proof}

The hypothesis fails exactly when some nontrivial $\rho$ has a direction that the
contractions $m_i$ asymptotically leave alone --- an $H$-fixed vector. The criterion is
easy to check in examples, and it is what separates the lifts.

\begin{theorem}[The simplex depends on the lift]\label{XXVII:thm:lift}
Let $G=S_3$, and let the class $C_i$ be the class of transpositions for every $i$.
\begin{enumerate}
\item For the constant lift $u_i=(12)$, $H=\{e,(12)\}$ and $\cA_c$ has exactly three
extremal $\mathrm{KMS}_\beta$ states, permuted transitively by $S_3$ with stabilizers of
order $2$.
\item For the alternating lift $u_{2k-1}=(12)$, $u_{2k}=(13)$, $H=S_3$ and $\cA_c$ has a
unique $\mathrm{KMS}_\beta$ state.
\end{enumerate}
In both cases $N_\beta=S_3$ and $\Prob(G/N_\beta)$ is a point: the Tannakian prediction is
right for the second lift and wrong for the first.
\end{theorem}

\begin{proof}
(1) The lifts commute, $c$ has constant increments $(12)$, and $H$ is the Mackey range of
the constant cocycle into the abelian group $A=\{e,(12)\}$, namely the annihilator of the
characters $\chi$ of $A$ with $\sum_it_i|1-\chi((12))|^2<\infty$; the nontrivial character
has $|1-\chi((12))|^2=4$ and $\sum t_i=\infty$, so only the trivial character qualifies and
$H=A$. Theorem~\ref{XXVII:thm:simplex} gives $\Prob(A\backslash S_3)$, three points.

(2) We verify the hypothesis of Theorem~\ref{XXVII:thm:mixing}. For the sign character,
$m_i=(1-t_i)/(1+t_i)$ and $\prod_{i>n}m_i=0$ since $\sum t_i=\infty$. For the two-dimensional
irreducible $\rho$, $\rho(u_i)$ is a reflection with fixed line $\ell_i$, $\ell_{2k-1}$ and
$\ell_{2k}$ making an angle of $60^\circ$; by \eqref{XXVII:eq:mean}, $m_i=Q_i+\lambda_iQ_i^\perp$
with $Q_i$ the projection onto $\ell_i$ and $\lambda_i=(1-t_i)/(1+t_i)$. For a unit vector $v$
either $\|Q_i^\perp v\|\ge\frac12$ or $\|Q_{i+1}^\perp v\|\ge\frac12$, the two lines being
$60^\circ$ apart. In the first case $\|m_iv\|^2\le1-\frac14(1-\lambda_i^2)$. In the second,
$\|m_iv-v\|\le1-\lambda_i\le2t_i\le\frac14$ once $t_i\le\frac18$, so
$\|Q_{i+1}^\perp m_iv\|\ge\frac14$ and $\|m_{i+1}m_iv\|^2\le1-\frac1{16}(1-\lambda_{i+1}^2)$.
As $1-\lambda_i^2\ge t_i$, in both cases
$\|m_{i+1}m_i\|_{\op}\le1-\frac1{32}\min(t_i,t_{i+1})$, so
\[
\|M_{(n,m]}\|_{\op}\le\prod_{k:\,n<2k-1<2k\le m}\Bigl(1-\tfrac1{32}\min(t_{2k-1},t_{2k})\Bigr)\le\exp\Bigl(-\tfrac1{32}\sum_{k}\min(t_{2k-1},t_{2k})\Bigr)\to0,
\]
because the $t_i$ are non-increasing, so $\min(t_{2k-1},t_{2k})=t_{2k}\ge t_{2k+1}$ and
$\sum_kt_{2k}\ge\frac12\sum_{i\ge2}t_i=\infty$. Theorem~\ref{XXVII:thm:mixing} applies.

For $N_\beta$: $\Xi_\beta(\mathrm{sgn},S)=4\sum t_i=\infty$ and
$\Xi_\beta(\rho,S)=\sum t_i\|I-\rho(u_i)\|_{\HS}^2=4\sum t_i=\infty$, so $\cT_\beta=\{\mathbf 1\}$
and $N_\beta=S_3$ for either lift, as it must, $N_\beta$ depending only on the classes.
\end{proof}

\begin{corollary}[No canonical non-abelian model]\label{XXVII:cor:nocanonical}
The $\mathrm{KMS}_\beta$ simplex of $\cA_c$ is not a function of the Frobenius classes
$(C_i)_i$: it depends on the lift, through the conjugacy class of the Mackey range $H$. The
lift-independent invariants of the construction are exactly those of $N_\beta=\langle\langle H\rangle\rangle$,
i.e.\ the Tannakian data of Chapter~\ref{ch:VI}.
\end{corollary}

\begin{remark}[Averaging over lifts does not help]
One may make the model canonical by adjoining the lift to the unit space: replace $G$ by
$G\times\prod_iG$ with Haar measure, let the $i$-th extra coordinate $k_i$ be a conjugator,
and use the lifts $k_iu_ik_i^{-1}$ in \eqref{XXVII:eq:c}. The extra coordinates are not moved by
any arrow, so the ergodic decomposition is fibred over the space of lifts, the KMS
simplex is the space of probability measures on the bundle of coset spaces
$H_k\backslash G$ over $k\in\prod G$, and the symmetry $G$ acts within the fibres: the
extreme points are not a single $G$-orbit. Canonicity is bought at the price of the
transitivity that defines symmetry breaking.
\end{remark}

\subsection{When the lift does not matter}

\begin{proposition}[Rigid Tannakian subgroups]\label{XXVII:prop:rigid}
If $N_\beta$ has no proper closed subgroup whose normal closure in $G$ is $N_\beta$ --- for
instance if $N_\beta$ is trivial, or of prime order, or more generally if every proper
closed subgroup of $N_\beta$ is contained in a proper closed $G$-normal subgroup of
$N_\beta$ --- then $H=N_\beta$ for every lift, and
$\mathrm{KMS}_\beta(\cA_c)\cong\Prob(G/N_\beta)$ with the transitive $G$-action, for every lift.
\end{proposition}

\begin{proof}
By Corollary~\ref{XXVII:cor:normalclosure}, $H\le N_\beta$ and $\langle\langle H\rangle\rangle=N_\beta$;
the hypothesis leaves only $H=N_\beta$.
\end{proof}

\begin{example}[A nontrivial $\Prob(G/N_\beta)$ realized for every lift]\label{XXVII:ex:A3}
Let $G=S_3$ and split $S=S_1\sqcup S_2$ with $\sum_{\pp\in S_1}\Ns\pp^{-\beta}=\infty$ and
$\sum_{\pp\in S_2}\Ns\pp^{-\beta}<\infty$; assign the class of $3$-cycles to the primes of $S_1$
and the class of transpositions to those of $S_2$. Then
$\Xi_\beta(\mathrm{sgn},S)=4\sum_{S_2}\Ns\pp^{-\beta}<\infty$, since the sign is trivial on
$3$-cycles, while $\Xi_\beta(\rho,S)\ge\sum_{S_1}\Ns\pp^{-\beta}\,\|I-\rho((123))\|_{\HS}^2=6\sum_{S_1}\Ns\pp^{-\beta}=\infty$
for the two-dimensional $\rho$. So $\cT_\beta=\{\mathbf 1,\mathrm{sgn}\}=\Rep(S_3/A_3)$ and
$N_\beta=A_3$, of prime order; by Proposition~\ref{XXVII:prop:rigid}, for every lift the ordered
model has exactly two extremal $\mathrm{KMS}_\beta$ states, exchanged by the transpositions
of $S_3$ and fixed by $A_3$: the simplex is $\Prob(S_3/A_3)$, a non-abelian symmetry group
breaking to a normal subgroup, as the Tannakian formalism predicts. Chebotarev's theorem,
used as in the constructions of [Ch.~\ref{ch:I}, \S5], supplies such splittings from an
$S_3$-extension of $K$: the primes whose Frobenius class consists of $3$-cycles have density
$\frac13$, so their Dirichlet sum diverges for $\beta\le1$, and for $S_2$ any set of
transposition primes with convergent sum will do.
\end{example}

\section{Verdict}

\[
\begin{array}{@{}l@{\ \ }l@{\ \ }l@{}}
\text{Route or statement} & \text{Outcome} & \text{Where}\\\hline
\text{B: }\cB\rtimes J_S,\ \cB\text{ noncommutative} & \text{no-go: KMS states are traces on }\cB & \ref{XXVII:thm:B}\\
\text{A: isotropy}=N_\beta & \text{no-go: gives }\Irr(N_\beta)\text{ and a fixed point} & \ref{XXVII:thm:isoNoGo}\\
\text{A: skew product by }c & \Prob(H\backslash G),\ H\text{ the Mackey range} & \ref{XXVII:thm:simplex}\\
\rho\circ c\ \text{coboundary}\iff\Xi_\beta(\rho,S)<\infty & \text{true, every lift} & \ref{XXVII:thm:kakutani}\\
N_\beta=\langle\langle H\rangle\rangle;\ \Prob(G/N_\beta)\iff H\trianglelefteq G & \text{true} & \ref{XXVII:cor:normalclosure}\\
\text{mixing means}\Rightarrow H=G & \text{true} & \ref{XXVII:thm:mixing}\\
\text{simplex independent of the lift} & \text{false }(S_3) & \ref{XXVII:thm:lift}\\
\Prob(G/N_\beta)\ \text{for every lift} & N_\beta\text{ rigid, e.g. }A_3 & \ref{XXVII:prop:rigid},\ \ref{XXVII:ex:A3}
\end{array}
\]

The decisive statements are these. There is no non-abelian Bost--Connes system in which
the representation theory of the coefficient algebra or the isotropy of the groupoid
produces the symmetry breaking: Theorems \ref{XXVII:thm:B} and \ref{XXVII:thm:isoNoGo} are no-go
theorems in the strict sense, for every $\cB$ and every $\cG$. There is a non-abelian
Bost--Connes groupoid, principal, \'etale and amenable, in which the Frobenius classes
drive the dynamics through a cocycle with non-commuting increments, and its equilibrium
theory is completely determined: the simplex is $\Prob(H\backslash G)$, the Frobenius
Dirichlet series decides which representations are coboundaries by the non-abelian
Kakutani theorem, and $N_\beta$ is the normal closure of $H$. What the Frobenius classes do
not determine is $H$ itself. Symmetry breaking to $G/N_\beta$ is realized by this groupoid
exactly when the lifts mix enough to make $H$ normal, always when $N_\beta$ is rigid, and,
when the lifts are chosen to commute as in Chapter~\ref{ch:VII}, exactly when the abelian
essential-value group $E_\beta$ happens to be normal. A canonical
non-abelian system, one whose equilibrium states depend on the Frobenius classes alone,
does not exist within this framework; the canonical object is $N_\beta$, and the passage
from $N_\beta$ to a simplex requires a choice that the arithmetic does not make.

